\scanpage{88}
un morphisme fini étale. Ainsi $\Pi_1(X,I)$ apparaît comme quotient de $\Pi_1(C\deleted{\unsure{X}},I)$, donc $\widehat\Pi_1(X,I)$ quotient de $\widehat\Pi_1(C,I)$, et où $\Pi_1(C,I)$ apparaît comme un revêtement fini de $\Pi_1(U_{0,3},I)$ --- plus \deleted{précisément} \inserted{concrètement}, prenant le saturé $I'=f^{-1}f(I)=\unsure{f}^{-1}(I_0)$ de $I$, $\Pi_1(U_{0,3},I)$ est un sous-groupoïde \inserted{plein} de $\Pi_1(U_{0,3};I')$, qui lui est un rev. fini du groupoïde $\Pi_1(U_{0,3};I_0)$ --- et idem pour les complétés profinis. Ainsi, l'opération de $\Gamma$ sur $\widehat\Pi_1(X)$ est connue quand on connaît l'opération sur $\widehat\Pi_1(C,\widetilde I)$, qui est un rev. fini de $\widehat\Pi_1(U_{0,3};I_0)$. Si on connaît l'action de $\Gamma$ sur $\widehat\Pi_1(U_{0,3};I_0)$ lui-même, \deleted{\unsure{La} \unsure{on} doit pouvoir \inserted{\unsure{l'action|l'opération} sur $\widehat\Pi_1(U_{0,3})$}}
\margin{87}
